\section{Huber location: complete known-scale proofs}
\label{sec:annex-p15-huber-robust-location}

\begin{annexitem}{annex:p15.convex.location.minimizers}{thm:p15-convex-location-minimizers}{Minimizers, measurable selection and one-sided derivatives}
\begin{proof}
The finite sum \(Q_n\) is continuous, convex and coercive, so its
minimum is attained in a bounded closed interval. Its minimizing
set is closed and convex, hence a nonempty compact interval.
Strict convexity rules out two distinct minimizers by evaluating
at their midpoint.

Joint continuity of \(Q_n(x,t)\) makes
\(m(x)=\inf_{q\in\mathbb Q}Q_n(x,q)\) measurable.
For rational \(r\), the restricted infima over rational \(q\leq r\)
and \(q\geq r\) are also measurable; by continuity and coercivity
they are the actual minima over those closed half-lines.
If \(l(x)\) is the left endpoint of the minimizing interval, then
\[
 \{l<a\}=\bigcup_{\substack{r\in\mathbb Q\\r<a}}
              \{\inf_{t\leq r}Q_n(x,t)=m(x)\}.
\]
This proves measurability of \(l\); the corresponding right-half-line
formula proves measurability of the right endpoint. Their midpoint
is measurable. The identity
\(Q_n(x+a,t+a)=Q_n(x,t)\) translates both endpoints, proving equivariance.

The right and left derivatives with respect to the location are
\[
 Q'_{n,+}(t)=-\sum_i\psi_-(x_i-t),\qquad
 Q'_{n,-}(t)=-\sum_i\psi_+(x_i-t).
\]
A convex function is minimized exactly where its left derivative is
nonpositive and its right derivative is nonnegative: the supporting
slopes then make every displacement nondecreasing in value; conversely
a strictly negative right derivative or positive left derivative
permits a decrease. This gives the stated score conditions.
If the midpoint score \(U_n(t)<0\), the average of these derivatives
is positive, so the right derivative is positive and no minimizer
can lie strictly to the right. If a minimizer lies strictly left
of \(t\), monotonicity of the one-sided derivatives makes both
derivatives at \(t\) nonnegative, giving \(U_n(t)\leq0\).
The reflected implications also hold.
\end{proof}
\end{annexitem}

\begin{annexitem}{annex:p15.convex.location.consistency}{thm:p15-convex-location-consistency}{Extended population scores and the strong law}
\begin{proof}
For a coercive convex criterion, its nondecreasing derivative is
strictly negative sufficiently far left and strictly positive
sufficiently far right. Otherwise the criterion could not tend
to infinity in the respective direction.
Split \(\psi=\psi^+-\psi^-\).
At \(t_0\) both expectations are finite.
For \(t\geq t_0\), \(\psi^+(X-t)\leq\psi^+(X-t_0)\);
for \(t\leq t_0\), \(\psi^-(X-t)\leq\psi^-(X-t_0)\).
Thus at least one part is finite everywhere, making \(\lambda\)
well defined and nonincreasing.
As \(t\to\infty\), the positive-part expectation tends to zero by
dominated convergence. An event \(X\leq M\) of positive
probability eventually contributes a fixed positive amount to the
negative part. Hence \(\lambda(t)<0\) eventually.
Reflection proves the other sign.

For a fixed threshold, the usual integrable strong law applies if
both parts of the score have finite expectation.
If its positive expectation is infinite and its negative part integrable,
apply \cref{thm:p15-etemadi-strong-law} to each bounded truncation
of the positive part and to the negative part.
On the common probability-one event for integer truncation levels,
the sample-average lower limit exceeds every finite number.
Thus the average tends to \(+\infty\). The negative-infinite case
follows by reflection.

For each rational \(\delta>0\), the empirical score is eventually
positive at \(c-\delta\) and negative at \(c+\delta\).
The convex derivative implications then place every minimizer in
\([c-\delta,c+\delta]\). Intersect over a countable sequence
\(\delta\downarrow0\) to obtain almost-sure convergence for every
measurable minimizing choice. Almost-sure convergence implies
convergence in probability, since the indicators of a fixed positive
deviation tend to zero almost surely and are bounded by one.
\end{proof}
\end{annexitem}

\begin{annexitem}{annex:p15.convex.location.normality}{thm:p15-convex-location-normality}{The source's triangular-array argument with corrected threshold sandwiches}
\begin{proof}
Translate \(c\) to zero. First suppose \(B>0\), and fix real \(z\).
Put \(t_n=z/\sqrt n\),
\(\lambda_n=\lambda(t_n)\), and
\(v_n=E\psi(X-t_n)^2-\lambda_n^2\to B\).
For all sufficiently large \(n\), \(t_n\) lies in the moment
neighborhood and \(v_n>0\). The centered independent row entries
\[
 W_{n,j}=\{\psi(X_j-t_n)-\lambda_n\}/\sqrt{nv_n}
\]
have total variance one.
Choose \(\eta>0\) with finite score second moments at \(\pm\eta\).
For all sufficiently large \(n\), monotonicity gives
\[
 |\psi(X-t_n)|\leq
 u(X):=\max(|\psi(X-\eta)|,|\psi(X+\eta)|),\qquad Eu^2<\infty.
\]
The means \(\lambda_n\) are bounded, and \(v_n\) is bounded away
from zero. For any fixed positive Lindeberg cutoff, the row's
large-value event therefore forces \(u(X)\) beyond a deterministic
threshold tending to infinity. Its squared centered score is bounded
by \(2u^2+2\sup_n\lambda_n^2\).
Dominated convergence on those tail events proves the Lindeberg
condition. Apply \cref{thm:p15-clt-lindeberg-feller} to get a
standard normal row sum.
Since \(\sqrt n\lambda_n\to-az\), this gives
\[
 \Pr\{U_n(t_n)<0\},\quad\Pr\{U_n(t_n)\leq0\}
                   \longrightarrow\Phi(az/\sqrt B).
\]
Both strict and weak events have the same limit because the limiting
normal law has no atom at zero.

For any minimizing choice and any \(\delta>0\), the exact convex
implications give the valid sandwich
\[
 \Pr\{U_n(z/\sqrt n)<0\}
 \leq\Pr\{\sqrt n T_n\leq z\}
 \leq\Pr\{U_n((z+\delta)/\sqrt n)\leq0\}.
\]
The upper implication uses \(T_n<(z+\delta)/\sqrt n\), not an
assumed equality of the score and zero at a nonsmooth minimizer.
Take limits and then \(\delta\downarrow0\).
This proves the normal limit, including flat minimizing intervals
and score jumps.
If the two population derivatives differ, the same row argument
uses \(a_-\) for negative \(z\) and \(a_+\) for positive \(z\).
The resulting CDF is continuous at zero, with value \(1/2\);
the same sandwich proves the stated joined-half-normal law.

If \(B=0\), for any fixed \(z\ne0\) the score sum divided by
\(\sqrt n\) has variance \(v_n\to0\) and mean tending to \(-az\).
Chebyshev's inequality then makes its sign correct with probability
tending to one. The convex implications at positive and negative
\(z/\sqrt n\) show \(\sqrt nT_n\to0\) in probability.
\end{proof}
\end{annexitem}

\begin{annexitem}{annex:p15.nonconvex.location.qualification}{thm:p15-nonconvex-location-qualification}{A measurable derivative average and the consistency premise}
\begin{proof}
Set
\[
 A_n=\int_0^1\frac1n\sum_i\psi'(X_i-sT_n)\,ds .
\]
This is measurable and finite for every sample.
Uniform continuity gives a modulus
\(\omega(h)=\sup_{|u-v|\leq h}|\psi'(u)-\psi'(v)|\to0\)
as \(h\downarrow0\), with finite values at finite \(h\) by subdivision.
Thus
\(\left|A_n-n^{-1}\sum_i\psi'(X_i)\right|\leq\omega(|T_n|)\to0\)
in probability. The integrable strong law gives \(A_n\to a>0\).
The fundamental theorem of calculus applied to each score gives
the exact equation \(n^{-1}\sum_i\psi(X_i)=T_nA_n\).
The right side tends to zero in probability, while the left tends
to \(E\psi(X)\), proving that mean is zero.
The i.i.d. central limit theorem and
\cref{thm:p15-clt-perturbation} now give
\(\sqrt nT_n=(n^{-1/2}\sum_i\psi(X_i))/A_n
 \Rightarrow N(0,E\psi^2/a^2)\).
The event \(A_n=0\) has probability tending to zero and does not
affect this conclusion.
For the counterexample the empirical criterion is
\(n(t^2-1)^2\), whose minimum zero occurs exactly at \(\pm1\).
The displayed alternating minimizing choice cannot converge in
probability to a constant.
\end{proof}
\end{annexitem}

\begin{annexitem}{annex:p15.huber.location.computation}{thm:p15-huber-location-computation}{Piecewise linear scores and the modified sample}
\begin{proof}
Each summand is continuous and nonincreasing in \(t\).
Below the sample minimum its sum is positive, and above the maximum
it is negative, unless all values coincide at the unique zero.
Between consecutive breakpoints, observations below the active
interval contribute \(-k\), those above it \(k\), and the active
ones \(x_i-t\). This gives the exact affine formula.
A nonconstant affine function has the stated sole zero; a constant
piece is zero exactly when its two saturated counts agree.
At breakpoints the sum is continuous, so including zero endpoints
completes the finite enumeration. The outer rays are nonzero.
Convexity already proved that the collected zeros form an interval.

At a zero \(t_0\), the modified observation is
\(y_i=t_0+\psi_k(x_i-t_0)\). Thus
\(\sum_i(y_i-t_0)=0\), proving that its mean is \(t_0\).
Its residual scores at \(t_0\) are unchanged by clipping, so it
also minimizes the modified criterion. This preserves a minimizing
location; it need not preserve every tie-breaking rule for every
modified sample.
For three zeros and \(M>4/3\), at \(t=1/3\) the three zero
observations have score \(-1/3\) and the last has score one.
The score therefore vanishes. In a neighborhood of this point
the three zero observations remain active and the last stays
saturated, so the sum is strictly decreasing there. Monotonicity
then excludes any other zero, proving uniqueness.
For \((-3,3)\), both scores are saturated and cancel exactly on
\([-2,2]\); outside that interval the sum has the corresponding
strict sign. The sample means and medians stated in the body
follow directly from their definitions.
\end{proof}
\end{annexitem}

\begin{annexitem}{annex:p15.huber.normal.variance}{thm:p15-huber-normal-variance}{Symmetry, endpoint atoms and the exact worst variance}
\begin{proof}
The bounded continuous score makes \(\lambda_F(t)\) continuous.
Its Gaussian component has derivative
\[
 -\{\Phi(t+k)-\Phi(t-k)\}<0,
\]
and its contamination component is nonincreasing.
Since \(1-\varepsilon>0\), the population score is strictly decreasing.
Its limits at negative and positive infinity are \(k,-k\), by
bounded convergence. Symmetry makes \(\lambda_F(0)=0\), hence it
is the unique root and the consistency theorem applies.

The difference quotient of \(\psi_k(x-t)\) in \(t\) is bounded
in absolute value by one. At zero its right derivative is minus
the indicator of \((-k,k]\), and its left derivative is minus
the indicator of \([-k,k)\). Dominated convergence gives these
integrals as the one-sided population derivatives.
Symmetry makes the two endpoint masses equal, so the derivatives
agree and have magnitude
\(A_F(k)=F((-k,k))+\tfrac12F(\{-k,k\})\).
This is at least \((1-\varepsilon)A(k)>0\).
The squared score is bounded and continuous in its shift, giving
the second-moment continuity required by the normality theorem.

Normal integration by parts, using \(\phi'(x)=-x\phi(x)\), gives
\[
 E_\Phi\min(Z^2,k^2)
 =\int_{-k}^k x^2\phi(x)\,dx+2k^2(1-\Phi(k))
 =A(k)-2k\phi(k)+2k^2(1-\Phi(k)).
\]
The contamination contributes at most \(k^2\) to this numerator,
and contributes a nonnegative amount to the denominator's derivative.
This proves the upper bound \(V^*(k,\varepsilon)\).
For symmetric \(H\) entirely on \(|x|>k\), the two contributions
are exactly \(k^2\) and zero, attaining the bound.
Known-scale transformation multiplies the centered estimator and
its normal fluctuation by \(\sigma\), giving the factor \(\sigma^2\).
\end{proof}
\end{annexitem}

\begin{annexitem}{annex:p15.huber.least.favorable.normal}{thm:p15-huber-least-favorable-normal}{The normal-exponential density and its information certificate}
\begin{proof}
Let \(C(k)=A(k)+2\phi(k)/k\). Differentiation gives
\(C'(k)=-2\phi(k)/k^2<0\).
Its limit at zero is infinity, and at infinity it is one.
Thus every \(1/(1-\varepsilon)>1\) has exactly one positive solution.
Integrating the central normal piece and its two exponential tails
gives \(\int f_*=(1-\varepsilon)C(k)=1\).
For \(|x|>k\),
\(\rho_k(x)=x^2/2-(|x|-k)^2/2\leq x^2/2\).
Hence \(f_*\geq(1-\varepsilon)\phi\) everywhere, with equality in
the central interval. The excess divided by \(\varepsilon\) is
a nonnegative symmetric density of mass one, supported outside
that interval. This proves membership in the original class.
The score \(\psi_k\) is continuous, so \(f_*\) is positive and \(C^1\),
with \(f_*'=-\psi_k f_*\).

The normalization equation gives
\[
 \frac{\varepsilon}{1-\varepsilon}
       =2\phi(k)/k-2(1-\Phi(k)).
\]
Multiply by \(k^2\) and combine with the formula for \(B(k)\) to get
\((1-\varepsilon)B(k)+\varepsilon k^2=(1-\varepsilon)A(k)\).
Thus \(I_*=E_{F_*}\psi_k^2=(1-\varepsilon)A(k)\), and the
worst-variance formula is \(1/I_*\).

For a competitor \(\psi\in L^2(F_*)\), change variables to write
\(\lambda_\psi(t)=\int\psi(y)f_*(y+t)\,dy\).
The Lipschitz bound on \(\rho_k\) gives
\(f_*(y+t)\leq e^{k|t|}f_*(y)\), and
\(|f_*'(y+t)|\leq k e^{k|t|}f_*(y)\).
Since \(E_{F_*}|\psi|<\infty\), dominated differentiation on a fixed
neighborhood of zero is justified. It yields
\(a_\psi=-\int\psi f_*'=E_{F_*}(\psi\psi_k)\).
Cauchy--Schwarz gives \(a_\psi^2\leq E\psi^2\,I_*\);
equality holds exactly for a nonzero proportional score when
the denominator is nonzero. This proves the score saddlepoint's
lower comparison.

Finally let \(f\) be an absolutely continuous density function in
the contamination class with finite \(I(f)=\int f'^2/f\).
Cauchy--Schwarz gives \(\int|f'|\leq\sqrt{I(f)}\).
Its limits at both infinities therefore exist and are zero, by
integrability of \(f\). Integration by parts with the bounded
absolutely continuous \(\psi_k\) gives
\[
 a_f=\int_{-k}^kf=-\int\psi_k f',\qquad
 b_f=\int\psi_k^2f,\qquad I(f)\geq a_f^2/b_f .
\]
Here \(a_f\geq I_*\) and \(b_f\leq I_*\) by the same contamination
bounds and the normalization identity above. Thus \(I(f)\geq I_*\).
Equality forces \(a_f=b_f=I_*\) and equality in Cauchy--Schwarz,
so \(f'=-\psi_k f\) almost everywhere.
The density is continuous and \(f\geq(1-\varepsilon)\phi\) almost
everywhere; continuity extends this lower bound everywhere, so \(f>0\).
On every finite interval, \(e^{\rho_k}f\) is absolutely continuous
with derivative zero almost everywhere, hence constant.
Normalization gives exactly \(f=f_*\). This proves uniqueness in
the stated density class.
\end{proof}
\end{annexitem}

\begin{annexitem}{annex:p15.location.integrated.score.bound}{thm:p15-location-integrated-score-bound}{An explicit shrinking-prior information proof}
\begin{proof}
Finite information implies
\(\int|f'|\leq(\int f'^2/f)^{1/2}<\infty\).
Thus \(f\) has zero limits at both infinities and \(\int f'=0\).
Let
\[
 w(u)=\cos^2(\pi u/2)\,{\bf1}_{\{|u|\leq1\}}.
\]
It is a \(C^1\) probability density, vanishes at the two endpoints,
and direct integration gives
\(\int w'^2/w=\pi^2\); the ratio is integrated only where \(w>0\).
For \(a>0\), use the location prior
\(\pi_a(\theta)=a^{-1}w(\theta/a)\), of information \(\pi^2/a^2\).
Clip the given estimate to \([-a,a]\), obtaining \(\widehat T\).
For every \(\theta\) in that interval,
\[
 |\widehat T-\theta|\leq|T_n-\theta|,\qquad
 |\widehat T-\theta|\leq2a .
\]

The joint density of \(\theta,x_1,\ldots,x_n\) is
\(r=\pi_a(\theta)\prod_i f(x_i-\theta)\).
On its positive set its location score is
\[
 H=\pi_a'(\theta)/\pi_a(\theta)
                -\sum_i f'(x_i-\theta)/f(x_i-\theta).
\]
Conditional means of the observation scores are \(\int f'=0\);
independence therefore gives
\(EH^2=\pi^2/a^2+nI(f)\).
This is finite. Moreover \(|\widehat T-\theta|\leq2a\), so
the product with \(H\) is absolutely integrable by Cauchy--Schwarz.
For fixed observations, integrate
\((\widehat T-\theta)\partial_\theta r\) over \([-a,a]\).
Integration by parts has zero boundary term because \(\pi_a\)
vanishes there and the finite positive \(C^1\) likelihood is
continuous on that finite interval. The result is \(\int r\,d\theta\).
The proved absolute integrability permits Fubini, giving
\(E[(\widehat T-\theta)H]=1\).
Cauchy--Schwarz now yields
\[
 E(\widehat T-\theta)^2\geq
                  \frac1{nI(f)+\pi^2/a^2}.
\]

Choose \(a=C/(2\sqrt n)\).
The two clipping inequalities give
\(n(\widehat T-\theta)^2\leq\min(n(T_n-\theta)^2,C^2)\).
Translation equivariance makes the expectation of the right side
independent of \(\theta\) in the location family, equal to
\(E_f\min(nT_n^2,C^2)\).
Average over the prior and substitute \(a\) into the preceding
lower bound. This is exactly
\(1/(I(f)+4\pi^2/C^2)\).
The calculation did not differentiate the estimator or assume its
untruncated mean squared error finite.
\end{proof}
\end{annexitem}

\begin{annexitem}{annex:p15.huber.location.minimax}{thm:p15-huber-location-minimax}{Keeping the three limiting operations in their stated order}
\begin{proof}
For any equivariant estimator sequence, apply the preceding
finite-\(n\), finite-\(C\) inequality at the single member \(F_*\).
First take \(\liminf_n\), then increase \(C\) to infinity.
This gives risk at \(F_*\) at least \(1/I_*\).
The outer supremum over \(F\) can only increase it.
Thus every competitor has \(\mathcal R_\varepsilon(T)\geq1/I_*\).

For the clipped estimator, fix \(F\in\mathcal C_\varepsilon\).
Its proved normal limit gives, for each fixed \(C\), convergence
of the bounded continuous loss \(\min(z^2,C^2)\) to its expectation
under \(N(0,V_k(F))\).
Increasing \(C\) then gives \(V_k(F)\) by monotone convergence.
Only after these limits take the supremum over \(F\), obtaining
\(V^*(k,\varepsilon)=1/I_*\), attained at \(F_*\).
This proves the asserted minimax identity without interchanging
a supremum with a limit.
For \(\varepsilon=0\), the normal density has information one,
the same lower bound applies, and the sample mean has limiting
variance one, proving the endpoint.

For \(k>0\), direct calculation gives
\[
 0\leq |x|-\rho_k(x)/k\leq k/2
\]
for every \(x\). Thus the normalized criteria converge uniformly
to the absolute-residual criterion, whose minimizing interval is
the median interval, as its one-sided derivative counts show.
For a finite sample all clipped minimizers lie in its fixed
minimum--maximum interval. Uniform convergence places their limit
points in the median interval.
If that interval is a point this proves convergence to the median.
If it has positive length (an even sample with distinct middle
values), for all sufficiently small \(k\) the clipped minimizing
interval is the middle gap with \(k\) removed from each end, so
its midpoint is the midpoint median. Literal \(k=0\), in contrast,
gives a constant criterion and no such minimizing rule.
\end{proof}
\end{annexitem}

\begin{annexitem}{annex:p15.huber.asymmetric.bias}{thm:p15-huber-asymmetric-bias}{The exact maximal population displacement}
\begin{proof}
The same bounded-convergence and Gaussian-derivative argument as
before makes every population score continuous and strictly decreasing,
with limits \(k,-k\). It therefore has a unique finite zero.
The consistency theorem applies at that zero.
The normal score \(m_k\) is odd and strictly decreasing, with
derivative \(-\{\Phi(c+k)-\Phi(c-k)\}\) and range \((-k,k)\).
At the root,
\[
 (1-\varepsilon)m_k(c)
       =-\varepsilon E_H\psi_k(X-c),\qquad
 |E_H\psi_k(X-c)|\leq k .
\]
For \(\varepsilon<1/2\) this places \(|c|\leq c_*\) as defined.
For \(H=\delta_M\), \(M>c_*+k\), the last expectation at \(c_*\)
is exactly \(k\), attaining the positive bound.
Reflection attains the negative one.
For \(\varepsilon\geq1/2\) and any proposed bound \(x>0\),
choose \(M>x+k\). Then
\((1-\varepsilon)m_k(x)+\varepsilon k>0\), since \(m_k(x)>-k\).
Its root lies above \(x\), proving unbounded maximal displacement.

Integrate \((z-c)\phi(z)\) on \([c-k,c+k]\) and the two saturated
tails to get the displayed exact formula for \(m_k(c)\).
Its odd smooth expansion at zero is
\(m_k(c)=-A(k)c+O(c^3)\).
For fixed \(k\), invert this strictly decreasing function near zero:
the root first obeys \(c_*=O(\varepsilon/(1-\varepsilon))\);
substitution then gives the stated cubic remainder.
This is an approximation to the exact root, not an unjustified
finite-\(\varepsilon\) upper bound.

Finally an asymmetric mixture does not by itself identify its
stipulated normal-component location. By
\cref{thm:p15-robust-overlap}, equal-fraction neighborhoods of two
normal locations intersect when their sup-event variation distance
is at most \(\varepsilon/(1-\varepsilon)\).
For normal locations separated by \(h\), their densities cross at
the midpoint and direct half-line integration gives distance
\(2\Phi(|h|/2)-1\), tending to zero with \(h\).
Thus for any positive fraction, sufficiently close distinct normal
locations can describe the same contaminated law.
\end{proof}
\end{annexitem}

\begin{annexitem}{annex:p15.huber.location.failure.cases}{thm:p15-huber-location-failure-cases}{Asymmetric atoms and finite-risk nonuniformity}
\begin{proof}
For the stated two-atom \(H\), the two score values are \(k,-k\),
so its mean score is zero. At \(+k\), a positive location shift
moves the residual inside the unclipped interval and contributes
derivative \(-1\); a negative shift leaves it saturated and
contributes zero. The atom at \(-M\) remains saturated for all
sufficiently small shifts. This gives precisely the two slopes
in the body. The squared score is bounded and shift-continuous,
so the two-half-normal theorem applies. Its two slopes differ,
excluding a single centered normal limit.

Under symmetric point contamination at \(\pm M\), the observation
variance is \(1-\varepsilon+\varepsilon M^2\).
The i.i.d. mean theorem gives that asymptotic variance for the
sample mean. For \(M>k\), the clipped-score worst-variance formula
has already been proved and is independent of \(M\).

For the finite-sample assertion, let \(T\) be any measurable estimate
in the convex hull and let \(T'\) be an independent copy.
Under the finite-\(M\) normal/point mixture, \(ET^2<\infty\), since
\(|T|^2\leq\max_i|X_i|^2\leq\sum_iX_i^2\).
With \(p=(\varepsilon/2)^n\), the event of all observations equal
to \(M\) forces \(T=M\) and has probability \(p\); similarly
all equal to \(-M\) forces \(T=-M\).
The two independent opposite events give
\[
 \operatorname{Var}(T)=\tfrac12E(T-T')^2
                        \geq4p^2M^2.
\]
Increase \(M\) to prove the infinite supremum for each fixed \(n\).
This explains exactly why the earlier limiting-risk criterion
cannot be replaced by that finite-sample supremum.
\end{proof}
\end{annexitem}

\begin{annexitem}{annex:p15.fisher.minimizer.nonuniqueness}{thm:p15-fisher-minimizer-nonuniqueness}{A complete disjoint-support counterexample}
\begin{proof}
The density \(f\) integrates to one, is continuously differentiable
after extension by zero, and has derivative
\(4\pi\sin(\pi x)\cos(\pi x)\) on \((0,1)\).
Its information, with zero integrand where both density and derivative
vanish, is
\[
 I(f)=\int_0^1 8\pi^2\cos^2(\pi x)\,dx=4\pi^2.
\]
The translated density \(g\) has the same information and disjoint
positive support. For \(0<t<1\), the information of the mixture
splits over those supports as
\[
 I((1-t)f+tg)=(1-t)I(f)+tI(g)=4\pi^2.
\]
The same equality holds at the endpoints.
The parameter interval is compact, and the integral of every
continuous compactly supported function against the mixture is
an affine continuous function of \(t\). Its image is therefore
compact for the vague topology and is convex by construction.
Every distinct member minimizes the same information value.
The normal-contamination uniqueness argument above has the additional
positive-density constraint and a specific equality certificate,
so this example does not contradict that proved result.
\end{proof}
\end{annexitem}
